% !TEX root = ../main.tex
\chapter{Materiales de reproducibilidad}
\label{anexo:materiales}

\section{Inventario}

La entrega reproducible deberá incluir:

\begin{itemize}
  \item corpus experimental SIGSAT-1.0, inventario, manifiesto y huellas de las 16 fuentes;
  \item catálogo congelado de 135 fragmentos y reporte de validación técnica;
  \item perfil y respuestas de definición confirmadas;
  \item guía de elaboración de RF, RNF e historias;
  \item línea base manual final con registro de tiempo;
  \item código y pruebas del prototipo;
  \item \emph{prompts} y parámetros congelados;
  \item registros de recuperación, puntos de control, artefactos, versiones, trazabilidad, validación y auditoría Jev cuando corresponda;
  \item rúbrica e instrucciones para expertos;
  \item calificaciones originales, acuerdo y adjudicaciones; y
  \item hojas o scripts empleados para calcular medidas.
\end{itemize}

\section{Convención de versiones}

Cada corrida evaluada deberá identificar como mínimo el corpus y sus huellas, la definición confirmada, la versión del código, la fecha UTC, el modelo LLM, el modelo de embeddings, método y pesos de recuperación, \(top\_k\), reranker, segmentación, temperatura, umbrales y versión de \emph{prompts}. Las credenciales de servicios externos nunca forman parte del material publicado.

\section{Declaración de uso de contenido sintético}

SIGSAT describe una organización y participantes ficticios. Sus correos, entrevistas, conversaciones, procedimientos, registros y cifras operativas no representan declaraciones de personas reales ni contienen información personal o confidencial. La única fuente no sintética es la Resolución SPDP-SPD-2025-0006-R, documento oficial público cuya procedencia y huella se conservan en los materiales de control. Esta distinción debe mantenerse visible en el corpus y en la tesis.
